\subsection{存在定理}

我们保留上一节的假设与记号。回顾，若 \(\alpha, \beta \in \mathrm{B}\) 且 \(\alpha \neq \beta\) ，则 \(n(\beta, \alpha) \leq 0\) ；而且，若 \(n(\beta, \alpha) = 0\) ，则 \(n(\alpha, \beta) = 0\) （第六章 §1 第 1 节 公式 (8)）。对 B 中两个不同元素组成的偶 \((\alpha, \beta)\) ，置

\[
x _ {\alpha \beta} = (\mathrm{ad} x _ {\alpha}) ^ {1 - n (\beta , \alpha)} x _ {\beta} y _ {\alpha \beta} = (\mathrm{ad} x _ {- \alpha}) ^ {1 - n (\beta , \alpha)} x _ {- \beta}.
\]

于是 \(x_{\alpha \beta} \in \mathfrak{a}_{+}, y_{\alpha \beta} \in \mathfrak{a}_{-}\) 。若用 \(\theta\) 表示 \(\mathfrak{a}\) 的典则自同构，则 \(\theta(x_{\alpha \beta}) = y_{\alpha \beta}\) 。

引理 6. 设 \(\alpha, \beta \in \mathrm{B}\) 且 \(\alpha \neq \beta\) 。则

\[
[ \mathfrak {a} _ {+}, y _ {\alpha \beta} ] = 0 \quad [ \mathfrak {a} _ {-}, x _ {\alpha \beta} ] = 0.
\]

第二个公式可由第一个公式利用自同构 \(\theta\) 得出。为证明第一个公式，只需证明 \([x_{\gamma}, y_{\alpha \beta}] = 0\) 对一切 \(\gamma \in B\) 成立。我们区分三种情形。

情形 1：\(\gamma \neq \alpha\) 且 \(\gamma \neq \beta\) 。在此情形，\(x_{\gamma}\) 与 \(x_{-\alpha}\) 及 \(x_{-\beta}\) 交换，从而与 \(y_{\alpha \beta}\) 交换。

情形 2：\(\gamma = \beta\) 。在此情形，\(x_{\gamma}\) 与 \(x_{-\alpha}\) 交换，故

\[
\begin{array}{r l} & {[ x _ {\gamma}, y _ {\alpha \beta} ] = (\mathrm{ad} x _ {- \alpha}) ^ {1 - n (\beta , \alpha)} [ x _ {\gamma}, x _ {- \beta} ]} \\ & {\qquad = - (\mathrm{ad} x _ {- \alpha}) ^ {1 - n (\beta , \alpha)} h _ {\beta} = - n (\alpha , \beta) (\mathrm{ad} x _ {- \alpha}) ^ {- n (\beta , \alpha)} x _ {- \alpha}.} \end{array}
\]

若 \(n(\beta, \alpha) < 0\) ，此表达式为零，因为 (ad \(x_{-\alpha}\) ). \(x_{-\alpha} = 0\) 。若 \(n(\beta, \alpha) = 0\) ，则 \(n(\alpha, \beta) = 0\) 。在两种情形下，都有 \([x_{\gamma}, y_{\alpha\beta}] = 0\) 。

情形 3：\(\gamma = \alpha\) 。在 \(\mathfrak{a}\) 的自同态代数中，

\[
[ - \operatorname{ad} h _ {\alpha}, \operatorname{ad} x _ {- \alpha} ] = 2 \operatorname{ad} x _ {- \alpha}
\]

且 {[}ad \(x_{\alpha}\) , ad \(x_{-\alpha}\) {]} = -ad \(h_{\alpha}\) ；于是由 §1 引理 1，

\[
[ \operatorname{ad} x _ {\alpha}, (\operatorname{ad} x _ {- \alpha}) ^ {1 - n (\beta , \alpha)} ] = (1 - n (\beta , \alpha)) (\operatorname{ad} x _ {- \alpha}) ^ {- n (\beta , \alpha)} (- \operatorname{ad} h _ {\alpha} - n (\beta , \alpha)).
\]

因此，

\[
\begin{array}{c} [ x _ {\gamma}, y _ {\alpha \beta} ] = [ \text {ad} x _ {\alpha}, (\text {ad} x _ {- \alpha}) ^ {1 - n (\beta , \alpha)} ] x _ {- \beta} + (\text {ad} x _ {- \alpha}) ^ {1 - n (\beta , \alpha)} (\text {ad} x _ {\alpha}) x _ {- \beta} \\ = - (1 - n (\beta , \alpha)) (\text {ad} x _ {- \alpha}) ^ {- n (\beta , \alpha)} (\text {ad} h _ {\alpha} + n (\beta , \alpha)) x _ {- \beta} \\ + (\text {ad} x _ {- \alpha}) ^ {1 - n (\beta , \alpha)} (\text {ad} x _ {\alpha}) x _ {- \beta}. \end{array}
\]

而 \([h_{\alpha},x_{-\beta}] + n(\beta ,\alpha)x_{-\beta} = 0\) 且 \([x_{\alpha},x_{-\beta}] = 0\) ，故 \([x_{\gamma},y_{\alpha \beta}] = 0\) 。

引理 7. \(\mathfrak{n}\) （\(\mathfrak{a}_{+}\) 中由诸 \(x_{\alpha \beta} (\alpha, \beta \in \mathrm{B}, \alpha \neq \beta)\) 生成的理想）是 \(\mathfrak{a}\) 的理想。\(\mathfrak{a}_{-}\) 中由诸 \(y_{\alpha \beta} (\alpha, \beta \in \mathrm{B}, \alpha \neq \beta)\) 生成的理想是 \(\mathfrak{a}\) 的理想，且等于 \(\theta(\mathfrak{n})\) 。

设 \(\mathfrak{n}' = \sum_{\alpha, \beta \in \mathrm{B}, \alpha \neq \beta} kx_{\alpha \beta}\) 。由于每个 \(x_{\alpha \beta}\) 是 \(\mathfrak{a}\) 中的齐次元素，\([\mathfrak{a}^0, \mathfrak{n}'] \subset \mathfrak{n}'\) （引理 5 与命题 2）。设 U （相应地 V ）是 \(\mathfrak{a}\) （相应地 \(\mathfrak{a}_+\) ）的包络代数，\(\sigma\) 是 U 在 \(\mathfrak{a}\) 上由 \(\mathfrak{a}\) 的伴随表示所定义的表示。\(\mathfrak{a}\) 的由 \(\mathfrak{n}'\) 生成的理想是 \(\sigma(\mathrm{U})\mathfrak{n}'\) 。而 \(\mathfrak{a} = \mathfrak{a}_+ + \mathfrak{a}^0 + \mathfrak{a}_-\) （命题 2），\(\sigma(\mathfrak{a}_-)\mathfrak{n}' = 0\) （引理 6），且由前所述 \(\sigma(\mathfrak{a}^0)\mathfrak{n}' \subset \mathfrak{n}'\) 。由 Poincaré-Birkhoff-Witt 定理，\(\sigma(\mathrm{U})\mathfrak{n}' = \sigma(\mathrm{V})\mathfrak{n}'\) ，这就证明了引理的第一个断言。由此推出，\(\theta(\mathfrak{a}_+) = \mathfrak{a}_-\) 中由诸 \(\theta(x_{\alpha \beta}) = y_{\alpha \beta}\) （ \(\alpha, \beta \in \mathrm{B}, \alpha \neq \beta\) ）生成的理想是理想 \(\theta(\mathfrak{n})\) （\(\mathfrak{a}\) 中的）。证毕。

a 的理想 \(\mathfrak{n} + \theta(\mathfrak{n})\) 是分次的，因为它由齐次元素生成。因此，李代数 \(\mathfrak{a}/(\mathfrak{n}+\theta(\mathfrak{n}))\) 是一个 Q-分次李代数；在本节其余部分，将它记作 \(g_{B}\) ，或简记作 g。由命题 2，若 \(g^{\mu} \neq 0\) ，则 \(\mu \in Q_{+}\) ，或 \(\mu \in Q_{-}\) ，或 \(\mu = 0\) 。用 \(X_{\alpha}\) （相应地 \(H_{\alpha}, X_{-\alpha}\) ）表示 \(x_{\alpha}\) （相应地 \(h_{\alpha}, x_{-\alpha}\) ）在 g 中的典则像。鉴于 a ，n 与 \(\theta(\mathfrak{n})\) 的定义，由此推出，g 是由生成族 \(((X_{\alpha}, H_{\alpha}, X_{-\alpha}))_{\alpha \in B}\) 和诸关系式所定义的李代数

\[
[ H _ {\alpha}, H _ {\beta} ] = 0\tag{16}
\]

\[
[ H _ {\alpha}, X _ {\beta} ] - n (\beta , \alpha) X _ {\beta} = 0\tag{17}
\]

\[
[ H _ {\alpha}, X _ {- \beta} ] + n (\beta , \alpha) X _ {- \beta} = 0\tag{18}
\]

\[
[ X _ {\alpha}, X _ {- \alpha} ] + H _ {\alpha} = 0\tag{19}
\]

\[
[ X _ {\alpha}, X _ {- \beta} ] = 0 (\alpha \neq \beta)\tag{20}
\]

\[
(\operatorname{ad} X _ {\alpha}) ^ {1 - n (\beta , \alpha)} X _ {\beta} = 0 \quad (\alpha \neq \beta)\tag{21}
\]

\[
(\operatorname{ad} X _ {- \alpha}) ^ {1 - n (\beta , \alpha)} X _ {- \beta} = 0 \quad (\alpha \neq \beta).\tag{22}
\]

设 \(z \in \mathfrak{g}\) ，\(\mu \in \mathrm{Q}\) 。则 \(z \in \mathfrak{g}^{\mu}\) 当且仅当 \([H_{\alpha}, z] = \langle \mu, \alpha^{\vee} \rangle z\) 对一切 \(\alpha \in \mathrm{B}\) 成立。这由引理 5 得出。

由于 \(\mathfrak{a}^0\cap (\mathfrak{n} + \theta (\mathfrak{n})) = 0\) ，从 \(\mathfrak{a}^0\) 到 \(\mathfrak{g}^0\) 的典则映射是同构。因此，\((H_{\alpha})_{\alpha \in \mathrm{B}}\) 是向量空间 \(\mathfrak{g}^0\) 的一个基。\(\mathfrak{g}^0\) （\(\mathfrak{g}_{\mathrm{B}}\) 的交换子代数）将记作 \(\mathfrak{h}_{\mathrm{B}}\) ，或简记作 \(\mathfrak{h}\) 。存在唯一的同构 \(\mu \mapsto \mu_{\mathrm{B}}\) （从 V 到 \(\mathfrak{h}^*\) ），使 \(\langle \mu_{\mathrm{B}},H_{\alpha}\rangle = \langle \mu ,\alpha^{\vee}\rangle\) 对一切 \(\mu \in \mathrm{V}\) 与一切 \(\alpha \in \mathrm{B}\) 成立。

a 的对合自同构 \(\theta\) 通过过渡到商定义了 g 的一个对合自同构，它同样记作 \(\theta\) 。我们有 \(\theta(X_{\alpha}) = X_{-\alpha}\) （对 \(\alpha \in B \cup (-B)\) ），以及 \(\theta(H_{\alpha}) = -H_{\alpha}\) 。

定理 1. 设 R 是一个既约根系，B 是 R 的一个基。设 g 是由生成族 \(((X_{\alpha}, H_{\alpha}, X_{-\alpha}))_{\alpha \in \mathrm{B}}\) 和关系式 (16) 至 (22) 所定义的李代数。设 \(h = \sum_{\alpha \in B} kH_{\alpha}\) 。则 \((\mathfrak{g}, \mathfrak{h})\) 是一个分裂半单李代数。同构 \(\mu \mapsto \mu_{B}\) （从 V 到 \(h^{*}\) ）把 R 映为 \((\mathfrak{g}, \mathfrak{h})\) 的根系。对一切 \(\mu \in R\) ，\(g^{\mu}\) 是相对于根 \(\mu\) 的特征子空间。

证明将循引理 8、9、10、11 的途径进行。

引理 8. 设 \(\alpha \in \mathrm{B} \cup (-\mathrm{B})\) 。则 \(\operatorname{ad} X_{\alpha}\) 是局部幂零的。 \footnote{一个自同态 \(u\) （向量空间 \(V\) 上的）称为局部幂零的（或殆幂零的），如果对每个 \(v \in V\) ，存在正整数 \(n\) 使 \(u^n(v) = 0\) （参见交换代数 第四章 §1 第 4 节 定义 2）。此时 \(\exp(u)\) 或 \(e^u\) 由公式 \(e^u(v) = \sum_{n \geq 0} (1/n!)u^n(v)\) 对一切 \(v \in V\) 定义。}

假设 \(\alpha \in \mathrm{B}\) 。设 \(\mathfrak{g}'\) 是这样的 \(z \in \mathfrak{g}\) 组成的集合：\((\operatorname{ad} X_{\alpha})^p z = 0\) 对充分大的 \(p\) 成立。由于 \(\operatorname{ad} X_{\alpha}\) 是 \(\mathfrak{g}\) 的导子，\(\mathfrak{g}'\) 是 \(\mathfrak{g}\) 的子代数。由 (21)，\(X_{\beta} \in \mathfrak{g}'\) 对一切 \(\beta \in \mathrm{B}\) 成立。由 (17)、(19)、(20)，\(H_{\beta} \in \mathfrak{g}'\) 与 \(X_{-\beta} \in \mathfrak{g}'\) 对一切 \(\beta \in \mathrm{B}\) 成立。于是 \(\mathfrak{g}' = \mathfrak{g}\) ，从而 \(\operatorname{ad} X_{\alpha}\) 是局部幂零的。由于 \(\operatorname{ad} X_{-\alpha} = \theta (\operatorname{ad} X_{\alpha})\theta^{-1}\) ，可见 \(\operatorname{ad} X_{-\alpha}\) 是局部幂零的。

我们将会看到 g 是有限维的，从而 \(ad X_{\alpha}\) 实际上是幂零的。

引理 9. 设 \(\mu, \nu \in \mathbf{Q}\) ，\(w \in \mathrm{W}(\mathrm{R})\) 满足 \(w\mu = \nu\) 。存在 \(\mathfrak{g}\) 的自同构把 \(\mathfrak{g}^{\mu}\) 变到 \(\mathfrak{g}^{\nu}\) 。

对一切 \(\alpha \in \mathrm{B}\) ，设 \(s_{\alpha}\) 是 V 的由 \(\alpha\) 定义的反射。由于 \(\mathrm{W}(\mathrm{R})\) 由诸 \(s_{\alpha}\) 生成（第六章 §1 第 5 节 注 1），只需就 \(w = s_{\alpha}\) 的情形证明引理。鉴于引理 8，我们可以定义

\[
\theta_ {\alpha} = e ^ {\mathrm{ad} X _ {\alpha}} e ^ {\mathrm{ad} X _ {- \alpha}} e ^ {\mathrm{ad} X _ {\alpha}}.
\]

与第一章 §6 第 8 节一样地验证，\(\theta_{\alpha}\) 是 \(\mathfrak{g}\) 的自同构。我们有

\[
\begin{array}{l} \theta_ {\alpha} (H _ {\alpha}) = e ^ {\mathrm{ad} X _ {\alpha}} e ^ {\mathrm{ad} X _ {- \alpha}} (H _ {\beta} - n (\alpha , \beta) X _ {\alpha}) \\ \qquad = e ^ {\mathrm{ad} X _ {\alpha}} \bigg (H _ {\beta} - n (\alpha , \beta) X _ {\alpha} + n (\alpha , \beta) X _ {- \alpha} - n (\alpha , \beta) H _ {\alpha} - \frac {n (\alpha , \beta)}{2} 2 X _ {- \alpha} \bigg) \\ \qquad = e ^ {\mathrm{ad} X _ {\alpha}} \left(H _ {\beta} - n (\alpha , \beta) H _ {\alpha} - n (\alpha , \beta) X _ {\alpha}\right) \\ \qquad = H _ {\beta} - n (\alpha , \beta) H _ {\alpha} - n (\alpha , \beta) X _ {\alpha} - n (\alpha , \beta) X _ {\alpha} - n (\alpha , \beta) (- 2 X _ {\alpha}) \\ \qquad = H _ {\beta} - n (\alpha , \beta) H _ {\alpha}. \end{array}
\]

若 \(z\in \mathfrak{g}^{\mu}\)

\[
\begin{array}{r l} & {[ H _ {\beta}, \theta_ {\alpha} ^ {- 1} z ] = \theta_ {\alpha} ^ {- 1} [ H _ {\beta} - n (\alpha , \beta) H _ {\alpha}, z ]} \\ & {\qquad = \theta_ {\alpha} ^ {- 1} (\langle \mu , \beta^ {\vee} \rangle z - n (\alpha , \beta) \langle \mu , \alpha^ {\vee} \rangle z)} \\ & {\qquad = \langle \mu - \langle \alpha^ {\vee}, \mu \rangle \alpha , \beta^ {\vee} \rangle \theta_ {\alpha} ^ {- 1} z = \langle s _ {\alpha} \mu , \beta^ {\vee} \rangle \theta_ {\alpha} ^ {- 1} z,} \end{array}
\]

故 \(\theta_{\alpha}^{-1}z\in\mathfrak{g}^{s_{\alpha}\mu}\) 。这表明 \(\theta_{\alpha}^{-1}\mathfrak{g}^{\mu}\subset\mathfrak{g}^{s_{\alpha}\mu}\) 。由于 \(\theta_{\alpha}\) 是自同构，且此包含关系对一切 \(\mu\in Q\) 成立，可见 \(\theta_{\alpha}^{-1}\mathfrak{g}^{\mu}=\mathfrak{g}^{s_{\alpha}\mu}\) ，这就证明了引理。

引理 10. 设 \(\mu\in Q\) ，并假设 \(\mu\) 不是根的倍数。存在 \(w\in\mathrm{W}(\mathrm{R})\) ，使 \(w\mu\) 关于基 B 的坐标中有一些 \(> 0\) ，另一些 \(< 0\)。

设 \(V_{R}\) 是向量空间 \(Q \otimes_{Z} R\) ，R 在其中是一个根系。由假设，存在 \(f \in V_{R}^{*}\) 使 \(\langle f, \alpha \rangle \neq 0\) 对一切 \(\alpha \in R\) 成立，且 \(\langle f, \mu \rangle = 0\) 。存在 \(R^{\vee}\) 的一个房 C 使 \(f \in C\) 。由第六章 §1 第 5 节 定理 2 (i)，存在 \(w \in W(R)\) 使 wf 属于对应于 B 的房，换言之，\(\langle wf, \alpha \rangle > 0\) 对一切 \(\alpha \in B\) 成立。记 \(w\mu = \sum_{\alpha \in B} t_{\alpha}\alpha\) 。则

\[
0 = \langle f, \mu \rangle = \langle w f, w \mu \rangle = \sum_ {\alpha \in B} t _ {\alpha} \langle w f, \alpha \rangle ,
\]

这证明某些 \(t_{\alpha}\) 为 \(> 0\) ，另一些为 \(< 0\) 。

引理 11. 设 \(\mu \in \mathrm{Q}\) 。若 \(\mu \notin \mathrm{R} \cup \{0\}\) ，则 \(\mathfrak{g}^{\mu} = 0\) 。若 \(\mu \in \mathrm{R}\) ，则 \(\dim \mathfrak{g}^{\mu} = 1\) 。 1) 若 \(\mu\) 不是 \(\mathrm{R}\) 中元素的倍数，则存在 \(w \in \mathrm{W}\) 使 \(w\mu \notin \mathrm{Q}_{+} \cup \mathrm{Q}_{-}\) （引理 10），于是 \(\mathfrak{a}^{w\mu} = 0\) ，\(\mathfrak{g}^{w\mu} = 0\) ，从而 \(\mathfrak{g}^{\mu} = 0\) （引理 9）。

\begin{enumerate}
\def\labelenumi{\arabic{enumi})}
\setcounter{enumi}{1}
\item
  设 \(\alpha \in \mathrm{B}\) ，\(m\) 是一个整数。由于 \(\mathfrak{a}_{+}\) 是以 \((x_{\alpha})_{\alpha \in \mathrm{B}}\) 为基础族的自由李代数，我们有 \(\dim \mathfrak{a}^{\alpha} = 1\) ，且 \(\mathfrak{a}^{m\alpha} = 0\) 对 \(m > 1\) 成立（第二章 §2 第 6 节 命题 4）。因此 \(\dim \mathfrak{g}^{\alpha} \leq 1\) ，且 \(\mathfrak{g}^{m\alpha} = 0\) 对 \(m > 1\) 成立。我们不可能有 \(\mathfrak{g}^{\alpha} = 0\) ，因为这将推出 \(x_{\alpha} \in \mathfrak{n} + \theta \mathfrak{n}\) ，从而 \(\mathfrak{n} + \theta \mathfrak{n}\) 包含 \(h_{\alpha} = -[x_{\alpha}, x_{-\alpha}]\) ，然而 \(\mathfrak{a}^0 \cap (\mathfrak{n} + \theta \mathfrak{n}) = 0\) 。于是，\(\dim \mathfrak{g}^{\alpha} = 1\) 。
\item
  若 \(\mu \in \mathrm{R}\) ，存在 \(w \in \mathrm{W}(\mathrm{R})\) 使 \(w(\mu) \in \mathrm{B}\) （第六章 §1 第 5 节 命题 15），故 \(\dim \mathfrak{g}^{\mu} = \dim \mathfrak{g}^{w\mu} = 1\) 。此外，若 \(n\) 是 \(>1\) 的整数，则 \(\mathfrak{g}^{nw(\mu)} = 0\) ，从而 \(\mathfrak{g}^{n\mu} = 0\) 。
\end{enumerate}

定理 1 的证明。

由于 \(\dim \mathfrak{g}^0 = \operatorname{CardB}\) ，由引理 11 推出 \(\mathfrak{g}\) 是有限维的，维数等于 \(\operatorname{CardB} + \operatorname{CardR}\) 。我们证明 \(\mathfrak{g}\) 是半单的。设 \(\mathfrak{k}\) 是 \(\mathfrak{g}\) 的一个交换理想。由于 \(\mathfrak{k}\) 在 \(\operatorname{ad}(\mathfrak{h})\) 之下稳定，\(\mathfrak{k} = (\mathfrak{k} \cap \mathfrak{h}) + \sum_{\mu \in \mathrm{R}} (\mathfrak{k} \cap \mathfrak{g}^\mu)\) 。

显然，对一切 \(\alpha \in \mathrm{B}\) ，\(\mathfrak{g}^{\alpha} + \mathfrak{g}^{-\alpha} + kH_{\alpha}\) 同构于 \(\mathfrak{sl}(2,k)\) 。鉴于引理 9，对一切 \(\mu \in \mathrm{R}\) ，\(\mathfrak{g}^{\mu}\) 包含于 \(\mathfrak{g}\) 的一个同构于 \(\mathfrak{sl}(2,k)\) 的子代数中；因此，\(\mathfrak{k} \cap \mathfrak{g}^{\mu} = 0\) ，故 \(\mathfrak{k} \subset \mathfrak{h}\) ；从而

\[
[ \mathfrak {k}, \mathfrak {g} ^ {\mu} ] \subset \mathfrak {k} \cap \mathfrak {g} ^ {\mu} = 0,
\]

故 \(\mu_{\mathrm{B}}(\mathfrak{k}) = 0\) 对一切 \(\mu \in \mathrm{R}\) 成立。由此推出 \(\mathfrak{k} = 0\) ，这就证明了 \(\mathfrak{g}\) 是半单的。

设 \(\mu \in \mathrm{R}\) 。存在 \(\alpha \in \mathrm{B}\) 使 \(\langle \mu, \alpha^{\vee} \rangle \neq 0\) ，于是 \((\operatorname{ad} H_{\alpha})|_{\mathfrak{g}^{\mu}}\) 是一个非零位似。因此，\(\mathfrak{h}\) 等于它在 \(\mathfrak{g}\) 中的正规化子，从而是 \(\mathfrak{g}\) 的一个嘉当子代数。对一切 \(u \in \mathfrak{h}\) ，\(\operatorname{ad} u\) 是可对角化的，故 \((\mathfrak{g}, \mathfrak{h})\) 是一个分裂半单李代数。

对一切 \(\mu \in \mathrm{R}\) ，显然 \(\mu_{\mathrm{B}}\) 是 \((\mathfrak{g},\mathfrak{h})\) 的根，且 \(\mathfrak{g}^{\mu}\) 是相应的特征子空间。\((\mathfrak{g},\mathfrak{h})\) 的根的数目是 \(\dim \mathfrak{g} - \dim \mathfrak{h} = \operatorname{Card}\mathrm{R}\) 。于是，映射 \(\mu \mapsto \mu_{\mathrm{B}}\) （从 V 到 \(\mathfrak{h}^*\) ）把 R 映为 \((\mathfrak{g},\mathfrak{h})\) 的根系。

